\section{Results}\label{sect:Results}

%%%%%%%%%%%%%%%%%%%%%%%%%%%%%%%%%%%%%%%%%%%
%   Introduce Results Section 
%%%%%%%%%%%%%%%%%%%%%%%%%%%%%%%%%%%%%%%%%%%

We present the numerical results of the lubrication-mediated quasi-potential model for both solid and deformable droplets, and begin by introducing a single test case which sits well within a stable parameter regime, using this to discuss the nuances of the model and how it compares to the DNS results in the deformable scenario. We proceed through a parameter sweep, highlighting the range of the model and general trends of droplet trajectories observed within a two-dimensional rebound case, all of which we present in the absence of gravity, as it is difficult for the droplet to detach under the effects of gravity in the two-dimensional system due to the stronger relative influence in the infinite-cylinder regime than a true spherical case \cite{vella2007surface}. Finally, we discuss the extensions to more complicated rebounds that can also be obtained, such as sustained periodic bouncing behaviour in the presence of Faraday forcing.

\begin{figure}
    \begin{center}
        \includegraphics[width = \textwidth]{Figures/Fig2.eps}
        \caption{Comparison of lubrication-mediated model for the solid and liquid drop cases, and DNS results for a deformable droplet, with initial undeformed radius $R_0 = 0.2$mm released with velocity $W_0=-0.2$ms$^{-1}$ from a centre of mass  height of $2R_0$. Shown are the centre of mass and poles of both droplets, and the free surface directly beneath the south pole. Simulations are run until the centre of mass returns to a height $R_0$.The liquid drop has the physical properties of water, and the solid impactor has the same density, with all other parameters indicated in Table~\ref{tab: Appendix}.}\label{fig: Rebound Comparison (t,z)}
    \end{center}
\end{figure}

%%%%%%%%%%%%%%%%%%%%%%%%%%%%%%%%%%%%%%%%%%%
%   Single Droplet Rebound
%%%%%%%%%%%%%%%%%%%%%%%%%%%%%%%%%%%%%%%%%%%
\subsection{LM model investigation}

Within Figures \ref{fig: Rebound Comparison (t,z)}-\ref{fig: h,Q,P nondimmed} we first consider a droplet with radius $R_0 = 0.2$ mm and initial velocity of $W_0 = -0.2$ms$^{-1}$, and compare the cases where the droplet is either solid or liquid to highlight the effect of droplet deformation on the rebound behaviour. In both cases we take the liquid in the bath to be water, and we take the droplet to be water or solid with the density of water. Parameter values are presented in Appendix~\ref{app: Table}. To consider solid droplets, the surface of the sphere is fixed and not permitted to deform setting $c_n=0$, hence $\eta_d(\theta,t) = R_0$. For further discussion of how viscous damping during impact affects rebound dynamics (including consideration of the soft solid limit), please see Appendix \ref{app: Visc}.

 

Figure~\ref{fig: Rebound Comparison (t,z)} presents the trajectories in the solid and deformable lubrication-mediated models, the latter also being contrasted to their counterpart DNS results. The lubrication-mediated rebound in the deformable case shows good agreement with DNS on the whole, though has a much smoother contact trajectory, which can be explained by the enhanced viscosity used during impact causing it to behave more closely to the solid drop during contact. Post contact, the droplet poles in both computations show oscillation generated by the deformation during contact, however the DNS oscillations are larger. The solid case has a slower exit velocity, due to the solid drop losing more kinetic energy to the bath, and not storing energy during impact as we shall see below. Henceforth, results are plotted nondimensionally using a capillary timescale $t_\sigma = \sqrt{\rho R_0^3/\sigma}$ and pressure $\sigma/R_0$ and with radius and initial velocity being used for typical length scale and velocity. 
We discuss the choice of enhanced viscous damping briefly in Appendix~\ref{app: Visc}.


%% Energy 
Figure~\ref{fig: Energy Plots (t,E) } demonstrates the transfer and decay of energy in the system in both the solid and deformable rebound cases. We can write the energy balance as 
\begin{equation*}
    E_{\text{tot}}(t) = K\!E_{\text{drop}} + E_{\text{wave}}  + E_{\text{osc}} - \Delta E_{\text{damp}},
\end{equation*}
where $E_{\text{tot}}$ denotes the total energy of the system, which at the start of impact is stored entirely in the kinetic energy of the droplet $ K\!E_{\text{drop}}$. $E_{\text{wave}}$ is the energy within the bath's wave field and $E_{\text{osc}}$ is the energy within the oscillations in the droplet. 

\begin{figure}
    \begin{center}
        \begin{subfigure}{0.49\linewidth}
            \includegraphics[width = \textwidth]{Figures/Fig3a.eps}
        \end{subfigure}
        \hspace*{\fill}
        \begin{subfigure}{0.49\linewidth}
            \includegraphics[width = \textwidth]{Figures/Fig3b.eps}
        \end{subfigure}
        \caption{Energy balance for the solid and deformable impactors for the lubrication-mediated model, with parameters as shown in Figure \ref{fig: Rebound Comparison (t,z)}, and in particular initial radius $R_0 = 0.2$mm, and initial velocity $W_0=-0.2$ms$^{-1}$. The energy in the system starts in the droplet kinetic energy $K\!E_{\text{drop}}$, and is transferred to the deformation and waves on the bath, and, in the deformable case, to the droplet oscillations during impact. The system energy is computed to be the sum of these terms, from which we can the loss of system energy due to the dissipation within the system.} \label{fig: Energy Plots (t,E) }
    \end{center}
\end{figure}

Once impact begins energy is transferred to the free surface motion and (for deformable drops) to droplet deformation. The effect of drop deformation on the overall dynamics is clear: deformation lessens the energy transferred into the waves of the bath, instead transferring energy to the oscillations within the droplet. Some of the droplet energy is returned in a rebounding effect and contributes to the upwards velocity of the droplet, increasing its kinetic energy relative to the solid case. This results in the higher exit velocity seen in the deformable case despite a lower final total energy. The effect of viscous damping is seen through the loss of total energy which is enhanced in the deformable drop case. The Ohnesorge number appears naturally as a decay constant, with a typical decay over a capillary timescale of $\exp{(-2\pi Oh)} \approx 0.95$, matching qualitatively the computational results.


\begin{figure}
    \begin{center}
        \begin{subfigure}{0.48\linewidth}
            \includegraphics[width = \textwidth]{Figures/Fig4a.eps}
            \subcaption{Fourier modes of the droplet free surface ($n=2,3,\ldots,31,32$). The $n=2$ mode experiences the greatest excitation and governs most of the droplet behaviour, with higher order modes experiencing less initial excitation and decaying more quickly. \\}
        \end{subfigure}
        \hspace*{\fill}
        \begin{subfigure}{0.48\linewidth}
            \includegraphics[width = \textwidth]{Figures/Fig4b.eps}
            \subcaption{Movement of the poles and equator from the rest state of the droplet. The south pole at $\theta =0$  experiences the largest disturbance as the bottom of the droplet flattens, with widening evident at the equator $\theta =\pi/2$. Oscillations are clear after the initial impact. \\}
        \end{subfigure}
        \hfill
        \begin{subfigure}{\linewidth}
            \includegraphics[width = \textwidth]{Figures/Fig4c.eps}
            \subcaption{Snapshots of the deformable droplet rebound from point of contact at $t=t_{\text{imp}}$ to detachment at $t=t_{\text{liftoff}}$, including maximum deformation at 4.19 and maximum depth at 4.80. 
            }
        \end{subfigure}
        \caption{Droplet deformation and droplet rebound characteristics for the deformable droplet lubrication-mediated model, as shown in Figure \ref{fig: Rebound Comparison (t,z)}, with initial radius $R_0 = 0.2$mm, and initial velocity $W_0=-0.2$ms$^{-1}$.}\label{fig: Droplet Deformations}
    \end{center}
\end{figure} 

%Drop Oscillations 

While the drop oscillation energy is small we may investigate the deformation with both the values of the Fourier coefficients of the drop's surface ($c_n$), and the overall shape profile of the droplet over time. Fig \ref{fig: Droplet Deformations}(a) illustrates normalised modes $c_n/R_0$, (where $n=2,\ldots, 32$) and it is clear that the second mode is the largest to be excited, and after detachment remains the main component of the oscillation. In contrast, for $n>6$ the modes have negligible effect on the overall behaviour post impact. They are however all excited at first contact, and we see this spike downwards as the southern half of the droplet essentially tries to `flatten' out as the lubrication layer initially spreads out. We see this in Figure~\ref{fig: Droplet Deformations}(b), where the south pole is deformed just before the north pole. Interestingly in this plot it is even possible to see the surface waves ripple around the droplet as they are excited at the equator and then at the north pole. Again, post impact we see the droplet continue to oscillate with period $t_\sigma$. 


\begin{figure}
    \begin{center}
        \begin{subfigure}{0.49\linewidth}
            \includegraphics[width = \textwidth]{Figures/Fig5a.eps}
        \end{subfigure}
        \begin{subfigure}{0.49\linewidth}
            \includegraphics[width = \textwidth]{Figures/Fig5b.eps}
        \end{subfigure}
        \caption{Dimensionless air layer height for a deformable droplet impact for the lubrication-mediated model (left) and DNS (right) as shown in Figure \ref{fig: Rebound Comparison (t,z)}, with initial radius $R_0 = 0.2$mm, and initial velocity $W_0=-0.2$ms$^{-1}$. The time axis in both figures is restricted to the approximate contact time, which begins from time $t/t_\sigma \approx 3$, to highlight impact behaviour. At the edges of the lubrication region and during the pinch before lift-off the dimensional film thicknesses may be as small as $ 0.5\;\mu$m.}\label{fig: Air Layer (x,t)}
    \end{center}
\end{figure}



% Pressure and Height within deformable 
A key distinction between the work presented here and the current literature is the dynamic inclusion of the lubrication layer  within the model. In Figure~\ref{fig: Air Layer (x,t)} we show a spatio-temporal plot of the lubrication layer height, over $\Omega_{l^*}(t)$. There are three main stages of rebound. The first stage happens quickly as the droplet ``impacts'' the bath and the lubrication region spreads out on timescales shorter than $t_\sigma$. Rapidly, the lubrication region  reaches its maximum extent and in the second stage begins to recede as the droplet peels away from the bath. During this stage, a narrow region forms near the outer edge where the layer depth is much thinner than at the centre of the lubrication region. If coalescence were permitted by the model, this would be the point 
of first contact that is described in the literature \cite{thoroddsen2003air}. We could expect to see contact when this depth is sufficiently small for atomic scale effects (e.g. Van der Waals forces) to become relevant. The third stage is pinchoff, which happens on a very short timescale and with negative pressure (suction) near the pinchoff point. 

\begin{figure}
    \begin{center}
        \begin{subfigure}{0.65\linewidth}
            \includegraphics[width = \textwidth]{Figures/Fig6a.eps} 
            \caption{}\label{fig: Pressure Surf}
        \end{subfigure}
        \begin{subfigure}{0.34\linewidth}
            \includegraphics[width = \textwidth]{Figures/Fig6b.eps}
            \caption{}\label{fig: Pressure Slices}
        \end{subfigure}
    \end{center}
    \caption{Dimensionless pressure $P^*$ within the air layer as determined by equation \eqref{eq: Numeric LubricationEquation} for the deformable lubrication-mediated rebound presented in Figure \ref{fig: Rebound Comparison (t,z)}, with initial radius $R_0 = 0.2$ mm and initial velocity $W_0 = -0.2$ms$^{-1}$, presented as (a) a spatio-temporal plot (with impact beginning from time $t/t_\sigma \approx 3$), and (b) at different times. The large spike of pressure at initial ``impact'' quickly spreads out and retains a flattened profile for most of contact. As the pressure diminishes in the later stages a region of low pressure forms on the boundary and propagates inwards, generating a suction effect. Eventually the suction region reaches the centre of the droplet, before detachment. The black lines in panel (a) correspond to the snapshots within figure~\ref{fig: h,Q,P nondimmed}. }\label{fig: Pressures}
\end{figure}




\begin{figure}
    \begin{center}
        \begin{subfigure}{0.49\linewidth}
            \includegraphics[width = \textwidth, left]{Figures/Fig7a.eps}\caption{$t/t_\sigma=2.98$}
        \end{subfigure}
        \begin{subfigure}{0.49\linewidth}
            \includegraphics[width = \textwidth, left]{Figures/Fig7b.eps}\caption{$t/t_\sigma=3.59$}
        \end{subfigure}
        \begin{subfigure}{0.49\linewidth}
            \includegraphics[width = \textwidth, right]{Figures/Fig7c.eps}\caption{$t/t_\sigma=5.41$}
        \end{subfigure}
        \begin{subfigure}{0.49\linewidth}
            \includegraphics[width = \textwidth, right]{Figures/Fig7d.eps}\caption{$t/t_\sigma=7.69$}
        \end{subfigure}
        \caption{Four representative snapshots of nondimensional layer height $h^*$, pressure $P^*$ and flux $Q^*$ within the lubrication region $\Omega_{l^*}(t)$ during the deformable lubrication-mediated rebound presented in Figure \ref{fig: Rebound Comparison (t,z)} with initial radius $R_0 = 0.2$mm and velocity $W_0 = -0.2$ms$^{-1}$. The times correspond to solid lines on Figure~\ref{fig: Pressure Surf}, chosen to highlight the change in behaviour between at distinct time regions during impact; initial contact (stage 1), widest point of contact, slow detachment (stage 2), and pinch-off (stage 3). Note that the axes change between each snapshot.}\label{fig: h,Q,P nondimmed}
    \end{center}
\end{figure}

The centre lubrication region in both the lubrication-mediated model and DNS presented is of near constant thickness of  $\mathcal{O}(1)\ \mu$m and pressure $\mathcal{O}(\sigma/R_0)$ for most of the impact, a quasi-static air cushion. This was experimentally observed in Tang \textit{et al} \cite{tang2019bouncing}, and corroborated in DNS computations \cite{galeano2021capillary, alventosa2023inertio, sprittles2024gas}. The region where the layer pinches at its boundaries lines up with the maxima (at the transition between stage 1 and stage 2) and minima (during stage 2) at the edge of the pressure profile. 

In Figure~\ref{fig: Pressures}(a) the horizontal bars across the plot correspond to the plots shown in Figure~\ref{fig: h,Q,P nondimmed}. There we show the nondimensional lubrication layer height $h/R_0$, pressure profile $P^*=P/(\sigma/R_0)$, and flux $Q^*$ against $x/R_0$, thus also indicating the fraction of the unperturbed radius they take up. These four times were selected as representative of impact stages. 

Figure~\ref{fig: h,Q,P nondimmed}(a) shows the profiles that the height, pressure, and flux during stage 1. The height of the layer is monotonically increasing, mimicking the shape of the droplet. There is a large spike in pressure, monotonically decreasing to atmospheric pressure, and a monotonically increasing flux as air is being forced out of the layer.

Figure~\ref{fig: h,Q,P nondimmed}(b) corresponds to the widest contact region, which occurs before the drop has reached maximum depth. The pressure is near constant (and hence flux is near zero) for the majority of the lubrication region with a very slight maximum near its edge and a smooth decay to atmospheric pressure at the outer region of the layer. The maximum is more pronounced in three-dimensions \cite{phillips2024lubrication}. The narrowing of the layer at its edge is now apparent.

Figure~\ref{fig: h,Q,P nondimmed}(c) corresponds to stage 2. The layer edge is propagating inwards towards the centre of the drop, the trapped air region contracts and the narrowing is more pronounced. There is a suction zone (negative pressure) and corresponding (negative) flux of outside air towards the outer boundary of the layer.

In the final stage 3 of impact (Figure~\ref{fig: h,Q,P nondimmed}(d)), the pinched regions on both sides of the drop have come together in the centre, resulting in a very thin layer there, and monotonically increasing away from the origin. The pressure is negative, delaying the detachment, and monotonically increasing. The flux corresponds to air filling the gap. Subsequently, the drop detaches and moves into flight away from the surface of the bath. 


%%%%%%%%%%%%%%%%%%%%%%%%%%%%%%%%%%%%%%%%%%%
%   Wider Parameter Sweep 
%%%%%%%%%%%%%%%%%%%%%%%%%%%%%%%%%%%%%%%%%%%
\subsection{Parameter study}

Through varying dimensionless grouping values $W\!e$ and $Oh$ in the absence of gravity, we construct an understanding of wider trends that can be observed. We present a series of cases within a region of applicability of the model, noting that the assumption that the deformations are small begin to break down at higher Weber numbers. 

We begin this parameter study by first highlighting the extremal values the investigation, in particular looking at distinct $W\!e$ and $Oh$ runs in Figure~\ref{fig: Deformation comparisons}. Scaled to highlight differences in deformation, the figure shows the four cases at the same nondimensional times. The first snapshot is taken at $t_\text{imp}/t_\sigma$, and the last at $t_\text{detachment}/t_\sigma$, where we define $t_\text{detachment}$ as the time when the drop and bath detach as the distance between them is larger than the threshhold $\varepsilon R_0$. The middle columns showcase intermediate deformations, with the penultimate column corresponding to maximum depth. The figures highlight that deformation of both of the liquid interfaces is more strongly linked to the Weber number than the Ohnesorge number, however the balance between viscous and surface tension effects still plays a part, as the bottom row at largest $W\!e$ and smallest $Oh$ of the four rows shows greatest deformation in the drop and bath. The Reynolds number $Re=\sqrt{W\!e}/Oh$ increases from row 1 to row 4 and is approximately equal for the two middle rows. 

\begin{figure}
    \begin{center}
        \includegraphics[width = \textwidth]{Figures/Fig8.eps}
        \caption{The evolution of four parameter regimes for the deformable LM model from point of impact to detachment. The top two runs correspond to a smaller droplet of radius $R_0=0.05$mm, and the lower two correspond to a larger droplet with radius $R_0=0.7$mm. The Weber number was varied by changing the initial velocity of the droplets respectively.}\label{fig: Deformation comparisons}
    \end{center}
\end{figure}
 In what follows, and to agree with prior studies, the time of impact $t_\text{imp}$ is defined as when the centre of mass of the droplet passes $z=R_0$ and $t_{\text{liftoff}}$ is when the centre of mass passes $z=R_0$ upwards again.  The contact time $t_c=t_{\text{liftoff}}-t_\text{imp}$ is the difference between these two times. This choice of time to determine contact will be different from when the lubrication layer `switches on' in the model, but this is done for a more accurate comparison to results in the wider literature. The other values we will consider are the maximum depth reached by the south pole of the droplet which we denote as $\delta$, and present as fraction of radius, and the coefficient of restitution, the negative of ratio of separation velocity at $t_\text{liftoff}$ to contact velocity at $t_\text{imp}$. Summarising:
\begin{equation*}
   \alpha = -\frac{W_{\text{out}}}{W_{\text{in}}}, \qquad \delta = \frac{(Z_{min}-R_0)}{R_0} ,\qquad \tau = \frac{t_c}{t_\sigma},
\end{equation*}
where $Z_{min}$ is the lowest point the centre of mass achieves during impact.
\begin{figure}
    \begin{center}
        \begin{subfigure}{0.49\linewidth}
            \caption{}
            \includegraphics[width = \textwidth]{Figures/Fig9a.eps}
        \end{subfigure}
        \begin{subfigure}{0.49\linewidth}
        \caption{}
            \includegraphics[width = \textwidth]{Figures/Fig9b.eps}
        \end{subfigure}
        \begin{subfigure}{0.49\linewidth}
        \caption{}
            \includegraphics[width = \textwidth]{Figures/Fig9c.eps}
        \end{subfigure}
        \begin{subfigure}{0.49\linewidth}
        \caption{}
            \includegraphics[width = \textwidth]{Figures/Fig9d.eps}
        \end{subfigure} 
        \begin{subfigure}{0.49\linewidth}
        \caption{}
            \includegraphics[width = \textwidth]{Figures/Fig9e.eps}
        \end{subfigure}
        \begin{subfigure}{0.49\linewidth}
        \caption{}
            \includegraphics[width = \textwidth]{Figures/Fig9f.eps}
        \end{subfigure}
        \caption{Comparisons of Coefficient of restitution $\alpha$, penetration depth $\delta$ and dimensionless contact time $\tau$, to Weber number (a,c,e) and Ohnesorge number (b,d,e) for a parameter sweep of values. The dimensionless parameter not plotted is indicated by the colour gradient on the opposite graph. Lubrication-mediated numerical results are presented by circles ($\bullet$), and triangles are used to indicate DNS results ($\triangle$) which correspond to a subset of the deformable parameter cases. The respective model subset is indicated by the circles being filled.}\label{fig: Comp We Oh}
    \end{center}
\end{figure}

\begin{figure}
    \begin{center}
        \begin{subfigure}{\linewidth}
            \includegraphics[width = \textwidth]{Figures/Fig10a.eps}
        \end{subfigure}
        \begin{subfigure}{\linewidth}
            \includegraphics[width = \textwidth]{Figures/Fig10b.eps}
        \end{subfigure}
        \begin{subfigure}{\linewidth}
            \includegraphics[width = \textwidth]{Figures/Fig10c.eps}
        \end{subfigure}
        \caption{Coefficient of restitution $\alpha$, penetration depth $\delta$ and dimensionless contact time $\tau$ plotted against Reynolds number $Re$. The colours correspond to varying Weber numbers $W\!e$, with connected lines corresponding to fixed Ohnesorge numbers $Oh$. Here we have included the solid model results by the squares ($\square$), and as before circles ($\bullet$) denote the deformable model results, with solid lines highlighting their DNS counterpart results shown with triangles ($\blacktriangle$).}\label{fig: Reynolds}
    \end{center}
\end{figure}

Figures~\ref{fig: Comp We Oh} and \ref{fig: Reynolds} display the three parameters of interest $\alpha$, $\delta$, $\tau$, against $W\!e$, $Oh$, $Re$. We include solid and drop simulations for the lubrication-mediated model, and drop simulations for DNS.

From Figure~\ref{fig: Rebound Comparison (t,z)} we observed that
for a fixed parameter choice the solid sphere rebounded less (had a lower coefficient of restitution) than its deformable counterpart. This is demonstrated again within Figure~\ref{fig: Reynolds}, where solid rebounds are characterised by smaller coefficient of restitution values. We note that within the study the solid rebound cases present little variance in $\alpha$, except for a dropoff at lower $Re$. 

For the  deformable droplet, as shown in Figure~\ref{fig: Comp We Oh}, the coefficient of restitution generally decreases with $Oh$ and $W\!e$ (with the other parameter fixed). Exception to this are the larger Ohnesorge cases (corresponding to smaller droplets) where there is a local maximum in $\alpha$ as a function of $W\!e$. This has been observed in DNS and kinematic match modelling of three-dimensional droplets \cite{alventosa2023inertio}. In the respective work  axisymmetric 3D simulations (without resolving a lubrication layer by using a kinematic match model) are compared to other studies, including experimental observations. Broadly, between the three-dimensional data therein and our two-dimensional computations, general trends are preserved despite the restitution coefficients in three dimensions being approximately a factor of two larger than in two dimensions. Our lubrication-mediated two-dimensional model results and trends corroborate far better with the two-dimensional DNS results herein, the only differences being that DNS results are monotonic in $W\!e$ and have slightly higher coefficients of restitution, particularly at lowest $Oh$ and intermediate $W\!e$. 

 Maximum depth results, as shown in Figure~\ref{fig: Comp We Oh}, indicate strong agreement on trends between the model and DNS. Although the lubrication-mediated model always reached a lower depth than the DNS results, particularly at higher $W\!e$, both models show that penetration depth is nearly independent of $Oh$ and has a strong positive correlation with $W\!e$. We note that a similar discrepancy between DNS and the kinematic match model was observed in Alventosa \textit{et al} \cite{alventosa2023inertio} and that the discrepancy fell within the experimental range. As shown in Figure~\ref{fig: Reynolds}, the difference in maximum depth between the solid and deformable lubrication-mediated results was small.

The contact time $\tau$ varies slowly with $W\!e$ and $Oh$ numbers, showing a general slow increase with both. The DNS and model results are consistent with each other, with only minor differences. In the small $W\!e$ limit, we observe a local minimum in $\tau$ for the larger $Oh$ values, correlating with the maximum of $\alpha$. There is an inverse relationship between contact time and coefficient of restitution: indeed a good portion of the contact time as defined here is post detachment when the droplet is still below the undisturbed free surface. As seen from Figure~\ref{fig: Droplet Deformations} in that case the actual impact takes about $5 t_\sigma$. The contact time is even longer for the solid impacts, as the impactors are taking even longer to reach this $z=R_0$ threshold criteria. 

Overall the agreement between the LM model and the DNS is highly encouraging, with key trends replicated successfully and explainable discrepancies owing to modelling assumptions or anticipated numerical effects on both sides. More specifically, DNS results can be summarised as exhibiting stronger pool deformation on impact, larger coefficients of restitution and a wider variation interval in contact times than their reduced-order model counterparts. Figure~\ref{fig:varyingviscosity} in Appendix~\ref{app: Visc} encapsulates some of the above features for a typical comparison case. The viscosity enhancement in the model provides a useful handle for navigating the deformable to non-deformable impactor property set, while addressing numerical challenges in terms of stability in its discretisation scheme. The lower the value of this parameter (down to a tractable threshold level), the closer the results become to the DNS benchmark liquid impact calculation, with pool deformation depth converging towards DNS-obtained results for the lowest enhanced viscosity values used in the model in the liquid impact case. The scenarios described by more pronounced pool deformation naturally lead to stronger bouncing, expressed in terms of larger coefficients of restitution, with any differences arguably becoming more prominent as $W\!e$ increases and fully nonlinear DNS calculations are expected to start diverging from the simplified model counterpart results. The behaviour of the impactor during the spreading and rebound stages in the two approaches is another point of differentiation, with more nonlinear behaviour in moderate to large deformation cases in the DNS datasets visible even as part of earlier air layer height dynamics studies, as shown for example in Figure~\ref{fig: Air Layer (x,t)} in terms of the radial extent of the contact region. More deformed droplets on impact in the DNS lead to larger contact radii and marginally thinner air layers, while the rebound phase is also slightly accelerated by comparison, correlating with the earlier observations on the small coefficient of restitution differences between the two employed methodologies.

We  postulate that the impact depth $\delta$ depends little on $Oh$ because the maximum depth is achieved early in the impact cycle. Thus most energy loss up to this stage is due to the transfer of energy to motion in the bath and not to viscous effects. One might expect $\delta \sim \sqrt{W\!e}$ from a model of a linear spring whose initial condition is proportional to the impact velocity $W_0$. The contact time and coefficient of restitution, however, depend more strongly on viscosity. When the data is plotted against $\sqrt{W\!e} \, Oh$ in Figure~\ref{fig: Collapsed} we observe a better description of the behaviour and collapse of the data. A justification for this scaling arises from postulating an exponential  decay rate for a linear spring model with decay rate proportional to $Oh$. This results in a linearly approximated viscous loss proportional to $Oh$, for small $Oh$ and finite time. For the contact times, we note that some high Weber number DNS data fall above lower Weber number data in this scaling. We conjecture that this is due to the increase in the (integer) number of oscillations of the droplet during impact in the DNS. These additional oscillations are not present in the LM model due to the additional damping effects introduced there.


\begin{figure}
    \begin{center}
        \begin{subfigure}{0.49\linewidth}
            \caption{}
            \includegraphics[width = \textwidth]{Figures/Fig11a.eps}
        \end{subfigure}
        \begin{subfigure}{0.49\linewidth}
        \caption{}
            \includegraphics[width = \textwidth]{Figures/Fig11b.eps}
        \end{subfigure}
        \caption{Collapsed data for coefficient of restitution $\alpha$ and dimensionless contact time $\tau$, both plotted against $\sqrt{W\!e}\,Oh$. The colour gradient corresponds to Weber number $W\!e$. The circles $(\bullet)$ denote deformable LM model results, and DNS data is shown using triangles $(\blacktriangle)$.}\label{fig: Collapsed}
    \end{center}
\end{figure}











